\documentclass[11pt]{amsart}
\usepackage[margin=1.1in]{geometry}
\usepackage{amssymb,mathtools,booktabs,multicol}
\usepackage[hidelinks]{hyperref}
\newcommand{\ntotal}{3525}
\newcommand{\ncerts}{2532}
\newcommand{\nclosed}{1954}
\newcommand{\nclosedone}{389}
\newcommand{\nclosedtwo}{1565}
\newcommand{\nimproved}{578}
\newcommand{\nchained}{3}
\newcommand{\nopen}{1571}
\newcommand{\countstable}{\begin{tabular}{rlrrrr}
\toprule
crossings & KnotInfo & knots & settled & lowered & open\\
\midrule
11 & $[1,\,2]$ & 14 & 0 & 0 & 14\\
11 & $[1,\,3]$ & 9 & 0 & 9 & 0\\
11 & $[2,\,3]$ & 36 & 30 & 0 & 6\\
12 & $[1,\,2]$ & 87 & 25 & 0 & 62\\
12 & $[1,\,3]$ & 87 & 12 & 69 & 6\\
12 & $[2,\,3]$ & 278 & 237 & 0 & 41\\
12 & $[2,\,4]$ & 3 & 2 & 0 & 1\\
13 & $[1,\,2]$ & 890 & 225 & 0 & 665\\
13 & $[1,\,3]$ & 677 & 126 & 497 & 54\\
13 & $[1,\,4]$ & 3 & 1 & 2 & 0\\
13 & $[1,\,5]$ & 1 & 0 & 0 & 1\\
13 & $[2,\,3]$ & 1437 & 1294 & 0 & 143\\
13 & $[2,\,4]$ & 3 & 2 & 1 & 0\\
\midrule
all & & 3525 & 1954 & 578 & 993\\
\bottomrule
\end{tabular}}
\newcommand{\listone}{12a\_133\\12a\_200\\12a\_212\\12a\_218\\12a\_313\\12a\_390\\12a\_457\\12a\_474\\12a\_611\\12a\_639\\12a\_698\\12a\_709\\12a\_750\\12a\_818\\12a\_878\\12a\_882\\12a\_895\\12a\_916\\12a\_1069\\12a\_1092\\12a\_1094\\12a\_1102\\12a\_1115\\12a\_1208\\12n\_32\\12n\_70\\12n\_147\\12n\_382\\12n\_390\\12n\_444\\12n\_495\\12n\_554\\12n\_555\\12n\_597\\12n\_632\\12n\_784\\12n\_831\\13a\_5\\13a\_6\\13a\_12\\13a\_129\\13a\_138\\13a\_156\\13a\_158\\13a\_166\\13a\_180\\13a\_232\\13a\_247\\13a\_253\\13a\_307\\13a\_338\\13a\_342\\13a\_343\\13a\_353\\13a\_365\\13a\_369\\13a\_375\\13a\_381\\13a\_392\\13a\_396\\13a\_410\\13a\_505\\13a\_518\\13a\_530\\13a\_538\\13a\_550\\13a\_591\\13a\_610\\13a\_619\\13a\_680\\13a\_712\\13a\_714\\13a\_721\\13a\_730\\13a\_732\\13a\_734\\13a\_777\\13a\_845\\13a\_906\\13a\_916\\13a\_925\\13a\_988\\13a\_995\\13a\_1027\\13a\_1068\\13a\_1072\\13a\_1077\\13a\_1103\\13a\_1106\\13a\_1122\\13a\_1134\\13a\_1163\\13a\_1169\\13a\_1176\\13a\_1186\\13a\_1193\\13a\_1230\\13a\_1231\\13a\_1237\\13a\_1250\\13a\_1277\\13a\_1316\\13a\_1324\\13a\_1362\\13a\_1380\\13a\_1421\\13a\_1436\\13a\_1438\\13a\_1454\\13a\_1511\\13a\_1514\\13a\_1532\\13a\_1598\\13a\_1608\\13a\_1609\\13a\_1624\\13a\_1648\\13a\_1656\\13a\_1666\\13a\_1671\\13a\_1712\\13a\_1719\\13a\_1722\\13a\_1727\\13a\_1756\\13a\_1758\\13a\_1862\\13a\_2094\\13a\_2166\\13a\_2218\\13a\_2308\\13a\_2323\\13a\_2328\\13a\_2339\\13a\_2389\\13a\_2411\\13a\_2414\\13a\_2502\\13a\_2505\\13a\_2560\\13a\_2595\\13a\_2645\\13a\_2788\\13a\_2828\\13a\_2870\\13a\_2925\\13a\_3012\\13a\_3028\\13a\_3039\\13a\_3046\\13a\_3057\\13a\_3058\\13a\_3167\\13a\_3181\\13a\_3188\\13a\_3190\\13a\_3242\\13a\_3258\\13a\_3266\\13a\_3274\\13a\_3275\\13a\_3279\\13a\_3302\\13a\_3306\\13a\_3333\\13a\_3343\\13a\_3355\\13a\_3356\\13a\_3403\\13a\_3408\\13a\_3415\\13a\_3418\\13a\_3462\\13a\_3463\\13a\_3475\\13a\_3477\\13a\_3494\\13a\_3505\\13a\_3515\\13a\_3522\\13a\_3556\\13a\_3579\\13a\_3598\\13a\_3609\\13a\_3611\\13a\_3618\\13a\_3653\\13a\_3657\\13a\_3692\\13a\_3741\\13a\_3758\\13a\_3764\\13a\_3772\\13a\_3776\\13a\_3786\\13a\_3787\\13a\_3805\\13a\_3807\\13a\_3811\\13a\_3905\\13a\_3928\\13a\_3955\\13a\_3993\\13a\_4077\\13a\_4103\\13a\_4152\\13a\_4195\\13a\_4198\\13a\_4211\\13a\_4220\\13a\_4247\\13a\_4268\\13a\_4365\\13a\_4367\\13a\_4398\\13a\_4403\\13a\_4412\\13a\_4419\\13a\_4437\\13a\_4443\\13a\_4475\\13a\_4487\\13a\_4493\\13a\_4513\\13a\_4524\\13a\_4544\\13a\_4577\\13a\_4591\\13a\_4603\\13a\_4629\\13a\_4647\\13a\_4648\\13a\_4696\\13a\_4697\\13a\_4705\\13a\_4744\\13a\_4746\\13a\_4794\\13a\_4848\\13n\_7\\13n\_30\\13n\_50\\13n\_54\\13n\_78\\13n\_149\\13n\_153\\13n\_157\\13n\_231\\13n\_290\\13n\_330\\13n\_341\\13n\_347\\13n\_360\\13n\_370\\13n\_406\\13n\_418\\13n\_428\\13n\_529\\13n\_530\\13n\_561\\13n\_563\\13n\_620\\13n\_671\\13n\_813\\13n\_816\\13n\_839\\13n\_852\\13n\_855\\13n\_892\\13n\_932\\13n\_988\\13n\_994\\13n\_999\\13n\_1057\\13n\_1061\\13n\_1062\\13n\_1067\\13n\_1114\\13n\_1201\\13n\_1243\\13n\_1275\\13n\_1302\\13n\_1328\\13n\_1419\\13n\_1490\\13n\_1492\\13n\_1494\\13n\_1582\\13n\_1585\\13n\_1609\\13n\_1628\\13n\_1638\\13n\_1678\\13n\_1688\\13n\_1693\\13n\_1726\\13n\_1789\\13n\_1932\\13n\_1935\\13n\_1992\\13n\_2022\\13n\_2080\\13n\_2193\\13n\_2265\\13n\_2383\\13n\_2404\\13n\_2406\\13n\_2453\\13n\_2503\\13n\_2610\\13n\_2629\\13n\_2694\\13n\_2706\\13n\_2741\\13n\_2765\\13n\_2812\\13n\_2834\\13n\_2837\\13n\_2851\\13n\_2857\\13n\_2926\\13n\_2942\\13n\_2966\\13n\_2977\\13n\_2985\\13n\_3108\\13n\_3190\\13n\_3198\\13n\_3236\\13n\_3311\\13n\_3344\\13n\_3351\\13n\_3452\\13n\_3526\\13n\_3641\\13n\_3667\\13n\_3695\\13n\_3740\\13n\_3745\\13n\_3752\\13n\_3771\\13n\_3899\\13n\_3902\\13n\_3916\\13n\_3919\\13n\_3921\\13n\_3923\\13n\_3933\\13n\_4023\\13n\_4030\\13n\_4107\\13n\_4115\\13n\_4152\\13n\_4155\\13n\_4162\\13n\_4193\\13n\_4202\\13n\_4247\\13n\_4300\\13n\_4305\\13n\_4313\\13n\_4333\\13n\_4334\\13n\_4351\\13n\_4354\\13n\_4380\\13n\_4457\\13n\_4477\\13n\_4509\\13n\_4515\\13n\_4521\\13n\_4614\\13n\_4632\\13n\_4642\\13n\_4666\\13n\_4709\\13n\_4746\\13n\_4870\\13n\_4896\\13n\_4940\\13n\_5006\\13n\_5009\\13n\_5012\\13n\_5015\\13n\_5041\\13n\_5089\\13n\_5095\\13n\_5104\\13n\_5108}
\newcommand{\listtwo}{11a\_17\\11a\_20\\11a\_49\\11a\_51\\11a\_53\\11a\_64\\11a\_89\\11a\_144\\11a\_161\\11a\_163\\11a\_178\\11a\_192\\11a\_214\\11a\_240\\11a\_241\\11a\_248\\11a\_274\\11a\_278\\11a\_292\\11a\_293\\11a\_299\\11a\_319\\11a\_329\\11a\_338\\11a\_339\\11a\_354\\11a\_355\\11a\_360\\11a\_362\\11a\_366\\12a\_9\\12a\_29\\12a\_35\\12a\_39\\12a\_43\\12a\_47\\12a\_49\\12a\_52\\12a\_55\\12a\_64\\12a\_66\\12a\_67\\12a\_76\\12a\_85\\12a\_88\\12a\_96\\12a\_97\\12a\_103\\12a\_113\\12a\_114\\12a\_117\\12a\_136\\12a\_137\\12a\_143\\12a\_146\\12a\_152\\12a\_167\\12a\_186\\12a\_188\\12a\_193\\12a\_224\\12a\_254\\12a\_264\\12a\_276\\12a\_289\\12a\_296\\12a\_297\\12a\_338\\12a\_347\\12a\_370\\12a\_371\\12a\_392\\12a\_428\\12a\_437\\12a\_449\\12a\_460\\12a\_461\\12a\_480\\12a\_507\\12a\_508\\12a\_514\\12a\_523\\12a\_533\\12a\_542\\12a\_546\\12a\_553\\12a\_556\\12a\_570\\12a\_575\\12a\_592\\12a\_597\\12a\_610\\12a\_616\\12a\_628\\12a\_633\\12a\_636\\12a\_641\\12a\_650\\12a\_651\\12a\_653\\12a\_655\\12a\_657\\12a\_661\\12a\_665\\12a\_674\\12a\_686\\12a\_689\\12a\_692\\12a\_702\\12a\_703\\12a\_706\\12a\_730\\12a\_738\\12a\_753\\12a\_762\\12a\_763\\12a\_767\\12a\_783\\12a\_787\\12a\_792\\12a\_795\\12a\_796\\12a\_806\\12a\_811\\12a\_812\\12a\_813\\12a\_815\\12a\_816\\12a\_817\\12a\_822\\12a\_828\\12a\_841\\12a\_844\\12a\_846\\12a\_869\\12a\_873\\12a\_876\\12a\_880\\12a\_881\\12a\_885\\12a\_890\\12a\_899\\12a\_905\\12a\_912\\12a\_950\\12a\_974\\12a\_986\\12a\_994\\12a\_996\\12a\_1001\\12a\_1002\\12a\_1006\\12a\_1009\\12a\_1013\\12a\_1021\\12a\_1025\\12a\_1030\\12a\_1031\\12a\_1035\\12a\_1036\\12a\_1037\\12a\_1044\\12a\_1049\\12a\_1052\\12a\_1058\\12a\_1059\\12a\_1079\\12a\_1080\\12a\_1084\\12a\_1088\\12a\_1091\\12a\_1097\\12a\_1099\\12a\_1103\\12a\_1107\\12a\_1112\\12a\_1114\\12a\_1116\\12a\_1120\\12a\_1122\\12a\_1137\\12a\_1143\\12a\_1152\\12a\_1164\\12a\_1172\\12a\_1177\\12a\_1184\\12a\_1193\\12a\_1197\\12a\_1200\\12a\_1222\\12a\_1223\\12a\_1229\\12a\_1232\\12a\_1240\\12a\_1244\\12a\_1248\\12a\_1252\\12a\_1255\\12a\_1265\\12a\_1270\\12a\_1274\\12n\_27\\12n\_33\\12n\_52\\12n\_67\\12n\_104\\12n\_107\\12n\_148\\12n\_153\\12n\_164\\12n\_169\\12n\_173\\12n\_174\\12n\_185\\12n\_186\\12n\_190\\12n\_211\\12n\_229\\12n\_243\\12n\_247\\12n\_260\\12n\_269\\12n\_289\\12n\_305\\12n\_317\\12n\_330\\12n\_361\\12n\_402\\12n\_432\\12n\_472\\12n\_503\\12n\_509\\12n\_517\\12n\_542\\12n\_557\\12n\_585\\12n\_603\\12n\_627\\12n\_638\\12n\_639\\12n\_671\\12n\_672\\12n\_675\\12n\_682\\12n\_688\\12n\_691\\12n\_692\\12n\_694\\12n\_707\\12n\_712\\12n\_717\\12n\_722\\12n\_734\\12n\_743\\12n\_746\\12n\_755\\12n\_788\\12n\_796\\12n\_803\\12n\_805\\12n\_818\\12n\_820\\12n\_827\\12n\_830\\12n\_863\\12n\_866\\12n\_874\\12n\_880\\13a\_4\\13a\_8\\13a\_16\\13a\_25\\13a\_29\\13a\_35\\13a\_51\\13a\_52\\13a\_56\\13a\_59\\13a\_65\\13a\_71\\13a\_82\\13a\_84\\13a\_85\\13a\_88\\13a\_89\\13a\_97\\13a\_98\\13a\_103\\13a\_112\\13a\_122\\13a\_134\\13a\_150\\13a\_151\\13a\_154\\13a\_157\\13a\_160\\13a\_162\\13a\_170\\13a\_171\\13a\_185\\13a\_201\\13a\_203\\13a\_209\\13a\_214\\13a\_217\\13a\_225\\13a\_231\\13a\_236\\13a\_251\\13a\_257\\13a\_273\\13a\_275\\13a\_284\\13a\_295\\13a\_304\\13a\_309\\13a\_314\\13a\_333\\13a\_334\\13a\_345\\13a\_355\\13a\_362\\13a\_366\\13a\_370\\13a\_374\\13a\_393\\13a\_402\\13a\_403\\13a\_406\\13a\_408\\13a\_423\\13a\_426\\13a\_429\\13a\_433\\13a\_437\\13a\_438\\13a\_443\\13a\_451\\13a\_468\\13a\_473\\13a\_474\\13a\_480\\13a\_481\\13a\_496\\13a\_497\\13a\_508\\13a\_513\\13a\_525\\13a\_535\\13a\_536\\13a\_539\\13a\_540\\13a\_542\\13a\_544\\13a\_545\\13a\_546\\13a\_551\\13a\_564\\13a\_569\\13a\_571\\13a\_573\\13a\_574\\13a\_584\\13a\_589\\13a\_593\\13a\_594\\13a\_598\\13a\_599\\13a\_603\\13a\_606\\13a\_618\\13a\_620\\13a\_621\\13a\_622\\13a\_623\\13a\_625\\13a\_629\\13a\_631\\13a\_636\\13a\_643\\13a\_644\\13a\_656\\13a\_658\\13a\_661\\13a\_669\\13a\_679\\13a\_682\\13a\_686\\13a\_716\\13a\_718\\13a\_744\\13a\_754\\13a\_768\\13a\_774\\13a\_784\\13a\_789\\13a\_808\\13a\_819\\13a\_826\\13a\_866\\13a\_867\\13a\_895\\13a\_900\\13a\_910\\13a\_918\\13a\_924\\13a\_931\\13a\_933\\13a\_955\\13a\_964\\13a\_966\\13a\_976\\13a\_980\\13a\_991\\13a\_1002\\13a\_1020\\13a\_1041\\13a\_1045\\13a\_1050\\13a\_1051\\13a\_1063\\13a\_1066\\13a\_1069\\13a\_1070\\13a\_1099\\13a\_1102\\13a\_1120\\13a\_1129\\13a\_1137\\13a\_1142\\13a\_1147\\13a\_1149\\13a\_1158\\13a\_1182\\13a\_1183\\13a\_1188\\13a\_1192\\13a\_1203\\13a\_1233\\13a\_1235\\13a\_1245\\13a\_1246\\13a\_1255\\13a\_1258\\13a\_1259\\13a\_1283\\13a\_1302\\13a\_1306\\13a\_1331\\13a\_1332\\13a\_1349\\13a\_1355\\13a\_1359\\13a\_1371\\13a\_1372\\13a\_1377\\13a\_1378\\13a\_1381\\13a\_1396\\13a\_1400\\13a\_1404\\13a\_1405\\13a\_1406\\13a\_1414\\13a\_1419\\13a\_1440\\13a\_1442\\13a\_1452\\13a\_1461\\13a\_1464\\13a\_1465\\13a\_1466\\13a\_1472\\13a\_1475\\13a\_1495\\13a\_1503\\13a\_1524\\13a\_1541\\13a\_1542\\13a\_1543\\13a\_1544\\13a\_1546\\13a\_1547\\13a\_1550\\13a\_1553\\13a\_1566\\13a\_1567\\13a\_1580\\13a\_1583\\13a\_1590\\13a\_1603\\13a\_1614\\13a\_1618\\13a\_1621\\13a\_1629\\13a\_1630\\13a\_1634\\13a\_1641\\13a\_1643\\13a\_1652\\13a\_1660\\13a\_1683\\13a\_1692\\13a\_1697\\13a\_1698\\13a\_1703\\13a\_1706\\13a\_1718\\13a\_1721\\13a\_1757\\13a\_1762\\13a\_1794\\13a\_1820\\13a\_1825\\13a\_1827\\13a\_1854\\13a\_1859\\13a\_1864\\13a\_1865\\13a\_1867\\13a\_1895\\13a\_1901\\13a\_1922\\13a\_1927\\13a\_1928\\13a\_1943\\13a\_1946\\13a\_1952\\13a\_1955\\13a\_1957\\13a\_1960\\13a\_1961\\13a\_1972\\13a\_1982\\13a\_1984\\13a\_1985\\13a\_1987\\13a\_1989\\13a\_1997\\13a\_1999\\13a\_2018\\13a\_2025\\13a\_2030\\13a\_2032\\13a\_2037\\13a\_2042\\13a\_2052\\13a\_2054\\13a\_2055\\13a\_2064\\13a\_2076\\13a\_2077\\13a\_2078\\13a\_2081\\13a\_2088\\13a\_2092\\13a\_2093\\13a\_2098\\13a\_2107\\13a\_2119\\13a\_2132\\13a\_2139\\13a\_2147\\13a\_2159\\13a\_2167\\13a\_2185\\13a\_2194\\13a\_2228\\13a\_2238\\13a\_2245\\13a\_2249\\13a\_2255\\13a\_2263\\13a\_2265\\13a\_2268\\13a\_2269\\13a\_2270\\13a\_2284\\13a\_2286\\13a\_2287\\13a\_2290\\13a\_2310\\13a\_2333\\13a\_2335\\13a\_2359\\13a\_2360\\13a\_2370\\13a\_2372\\13a\_2375\\13a\_2376\\13a\_2386\\13a\_2387\\13a\_2390\\13a\_2395\\13a\_2399\\13a\_2400\\13a\_2405\\13a\_2407\\13a\_2408\\13a\_2409\\13a\_2419\\13a\_2421\\13a\_2428\\13a\_2434\\13a\_2438\\13a\_2443\\13a\_2450\\13a\_2452\\13a\_2453\\13a\_2462\\13a\_2474\\13a\_2476\\13a\_2478\\13a\_2485\\13a\_2520\\13a\_2526\\13a\_2533\\13a\_2536\\13a\_2538\\13a\_2543\\13a\_2550\\13a\_2552\\13a\_2554\\13a\_2561\\13a\_2565\\13a\_2578\\13a\_2593\\13a\_2598\\13a\_2600\\13a\_2604\\13a\_2606\\13a\_2607\\13a\_2610\\13a\_2611\\13a\_2616\\13a\_2628\\13a\_2636\\13a\_2641\\13a\_2644\\13a\_2655\\13a\_2656\\13a\_2679\\13a\_2683\\13a\_2697\\13a\_2706\\13a\_2709\\13a\_2711\\13a\_2724\\13a\_2732\\13a\_2735\\13a\_2751\\13a\_2757\\13a\_2761\\13a\_2763\\13a\_2767\\13a\_2769\\13a\_2772\\13a\_2782\\13a\_2784\\13a\_2786\\13a\_2790\\13a\_2813\\13a\_2814\\13a\_2816\\13a\_2821\\13a\_2823\\13a\_2829\\13a\_2840\\13a\_2843\\13a\_2853\\13a\_2856\\13a\_2863\\13a\_2865\\13a\_2866\\13a\_2872\\13a\_2880\\13a\_2881\\13a\_2886\\13a\_2888\\13a\_2897\\13a\_2911\\13a\_2920\\13a\_2937\\13a\_2959\\13a\_2970\\13a\_2981\\13a\_2982\\13a\_2984\\13a\_2994\\13a\_2996\\13a\_3000\\13a\_3007\\13a\_3021\\13a\_3030\\13a\_3040\\13a\_3044\\13a\_3045\\13a\_3059\\13a\_3060\\13a\_3062\\13a\_3068\\13a\_3074\\13a\_3080\\13a\_3083\\13a\_3089\\13a\_3090\\13a\_3093\\13a\_3099\\13a\_3100\\13a\_3112\\13a\_3114\\13a\_3120\\13a\_3127\\13a\_3129\\13a\_3132\\13a\_3133\\13a\_3141\\13a\_3146\\13a\_3153\\13a\_3155\\13a\_3156\\13a\_3159\\13a\_3161\\13a\_3166\\13a\_3174\\13a\_3176\\13a\_3178\\13a\_3179\\13a\_3182\\13a\_3185\\13a\_3192\\13a\_3207\\13a\_3212\\13a\_3217\\13a\_3223\\13a\_3225\\13a\_3230\\13a\_3232\\13a\_3235\\13a\_3237\\13a\_3239\\13a\_3240\\13a\_3244\\13a\_3248\\13a\_3253\\13a\_3255\\13a\_3257\\13a\_3259\\13a\_3261\\13a\_3265\\13a\_3267\\13a\_3277\\13a\_3283\\13a\_3290\\13a\_3297\\13a\_3300\\13a\_3308\\13a\_3313\\13a\_3324\\13a\_3327\\13a\_3332\\13a\_3334\\13a\_3336\\13a\_3341\\13a\_3344\\13a\_3371\\13a\_3373\\13a\_3379\\13a\_3386\\13a\_3388\\13a\_3390\\13a\_3400\\13a\_3401\\13a\_3426\\13a\_3428\\13a\_3436\\13a\_3444\\13a\_3460\\13a\_3468\\13a\_3478\\13a\_3479\\13a\_3482\\13a\_3492\\13a\_3500\\13a\_3504\\13a\_3508\\13a\_3512\\13a\_3531\\13a\_3534\\13a\_3535\\13a\_3537\\13a\_3551\\13a\_3553\\13a\_3559\\13a\_3564\\13a\_3567\\13a\_3570\\13a\_3577\\13a\_3592\\13a\_3594\\13a\_3595\\13a\_3596\\13a\_3597\\13a\_3605\\13a\_3615\\13a\_3621\\13a\_3631\\13a\_3639\\13a\_3652\\13a\_3663\\13a\_3667\\13a\_3669\\13a\_3676\\13a\_3686\\13a\_3688\\13a\_3690\\13a\_3691\\13a\_3699\\13a\_3704\\13a\_3714\\13a\_3716\\13a\_3731\\13a\_3736\\13a\_3737\\13a\_3743\\13a\_3745\\13a\_3752\\13a\_3754\\13a\_3759\\13a\_3792\\13a\_3795\\13a\_3804\\13a\_3812\\13a\_3822\\13a\_3823\\13a\_3827\\13a\_3844\\13a\_3845\\13a\_3847\\13a\_3864\\13a\_3869\\13a\_3872\\13a\_3883\\13a\_3888\\13a\_3896\\13a\_3899\\13a\_3912\\13a\_3922\\13a\_3926\\13a\_3927\\13a\_3940\\13a\_3943\\13a\_3960\\13a\_3962\\13a\_3972\\13a\_3975\\13a\_3976\\13a\_3977\\13a\_3979\\13a\_3998\\13a\_3999\\13a\_4005\\13a\_4009\\13a\_4012\\13a\_4021\\13a\_4025\\13a\_4043\\13a\_4046\\13a\_4048\\13a\_4049\\13a\_4050\\13a\_4057\\13a\_4059\\13a\_4065\\13a\_4069\\13a\_4072\\13a\_4080\\13a\_4083\\13a\_4087\\13a\_4088\\13a\_4089\\13a\_4090\\13a\_4093\\13a\_4096\\13a\_4101\\13a\_4106\\13a\_4108\\13a\_4109\\13a\_4113\\13a\_4118\\13a\_4120\\13a\_4123\\13a\_4125\\13a\_4128\\13a\_4134\\13a\_4149\\13a\_4153\\13a\_4162\\13a\_4170\\13a\_4179\\13a\_4184\\13a\_4191\\13a\_4193\\13a\_4206\\13a\_4207\\13a\_4208\\13a\_4216\\13a\_4225\\13a\_4226\\13a\_4229\\13a\_4231\\13a\_4233\\13a\_4235\\13a\_4241\\13a\_4254\\13a\_4255\\13a\_4259\\13a\_4260\\13a\_4263\\13a\_4265\\13a\_4281\\13a\_4286\\13a\_4287\\13a\_4290\\13a\_4293\\13a\_4304\\13a\_4310\\13a\_4312\\13a\_4316\\13a\_4320\\13a\_4321\\13a\_4322\\13a\_4328\\13a\_4345\\13a\_4353\\13a\_4362\\13a\_4368\\13a\_4369\\13a\_4376\\13a\_4377\\13a\_4379\\13a\_4381\\13a\_4382\\13a\_4385\\13a\_4386\\13a\_4388\\13a\_4390\\13a\_4393\\13a\_4409\\13a\_4421\\13a\_4422\\13a\_4424\\13a\_4429\\13a\_4431\\13a\_4445\\13a\_4454\\13a\_4457\\13a\_4459\\13a\_4464\\13a\_4468\\13a\_4478\\13a\_4480\\13a\_4481\\13a\_4482\\13a\_4483\\13a\_4484\\13a\_4492\\13a\_4499\\13a\_4502\\13a\_4503\\13a\_4508\\13a\_4517\\13a\_4518\\13a\_4531\\13a\_4535\\13a\_4536\\13a\_4541\\13a\_4549\\13a\_4551\\13a\_4552\\13a\_4554\\13a\_4555\\13a\_4557\\13a\_4565\\13a\_4567\\13a\_4568\\13a\_4571\\13a\_4574\\13a\_4580\\13a\_4589\\13a\_4599\\13a\_4606\\13a\_4609\\13a\_4622\\13a\_4624\\13a\_4636\\13a\_4644\\13a\_4645\\13a\_4646\\13a\_4656\\13a\_4657\\13a\_4660\\13a\_4667\\13a\_4673\\13a\_4691\\13a\_4704\\13a\_4707\\13a\_4709\\13a\_4718\\13a\_4725\\13a\_4732\\13a\_4733\\13a\_4737\\13a\_4740\\13a\_4747\\13a\_4757\\13a\_4758\\13a\_4766\\13a\_4767\\13a\_4770\\13a\_4771\\13a\_4787\\13a\_4790\\13a\_4792\\13a\_4796\\13a\_4798\\13a\_4807\\13a\_4809\\13a\_4811\\13a\_4812\\13a\_4813\\13a\_4817\\13a\_4820\\13a\_4824\\13a\_4826\\13a\_4832\\13a\_4835\\13a\_4837\\13a\_4840\\13a\_4842\\13a\_4843\\13a\_4849\\13a\_4850\\13a\_4851\\13a\_4853\\13a\_4858\\13a\_4859\\13a\_4861\\13a\_4864\\13a\_4865\\13a\_4869\\13a\_4873\\13n\_6\\13n\_17\\13n\_25\\13n\_32\\13n\_41\\13n\_44\\13n\_45\\13n\_75\\13n\_82\\13n\_113\\13n\_131\\13n\_150\\13n\_154\\13n\_177\\13n\_186\\13n\_187\\13n\_204\\13n\_215\\13n\_222\\13n\_225\\13n\_241\\13n\_247\\13n\_265\\13n\_274\\13n\_285\\13n\_300\\13n\_312\\13n\_317\\13n\_321\\13n\_338\\13n\_339\\13n\_342\\13n\_356\\13n\_368\\13n\_378\\13n\_384\\13n\_393\\13n\_429\\13n\_437\\13n\_438\\13n\_458\\13n\_486\\13n\_495\\13n\_500\\13n\_503\\13n\_517\\13n\_519\\13n\_535\\13n\_546\\13n\_550\\13n\_569\\13n\_596\\13n\_601\\13n\_616\\13n\_618\\13n\_622\\13n\_629\\13n\_631\\13n\_663\\13n\_670\\13n\_675\\13n\_679\\13n\_680\\13n\_681\\13n\_691\\13n\_709\\13n\_723\\13n\_732\\13n\_744\\13n\_747\\13n\_755\\13n\_767\\13n\_771\\13n\_773\\13n\_776\\13n\_787\\13n\_832\\13n\_841\\13n\_858\\13n\_895\\13n\_900\\13n\_907\\13n\_914\\13n\_916\\13n\_929\\13n\_935\\13n\_943\\13n\_947\\13n\_961\\13n\_962\\13n\_965\\13n\_974\\13n\_981\\13n\_983\\13n\_1004\\13n\_1008\\13n\_1011\\13n\_1029\\13n\_1043\\13n\_1049\\13n\_1058\\13n\_1063\\13n\_1072\\13n\_1075\\13n\_1077\\13n\_1080\\13n\_1086\\13n\_1095\\13n\_1097\\13n\_1104\\13n\_1106\\13n\_1115\\13n\_1117\\13n\_1176\\13n\_1180\\13n\_1187\\13n\_1192\\13n\_1204\\13n\_1205\\13n\_1212\\13n\_1222\\13n\_1225\\13n\_1229\\13n\_1235\\13n\_1238\\13n\_1239\\13n\_1240\\13n\_1273\\13n\_1278\\13n\_1294\\13n\_1300\\13n\_1323\\13n\_1329\\13n\_1332\\13n\_1335\\13n\_1341\\13n\_1344\\13n\_1382\\13n\_1383\\13n\_1390\\13n\_1406\\13n\_1415\\13n\_1417\\13n\_1421\\13n\_1439\\13n\_1498\\13n\_1503\\13n\_1504\\13n\_1508\\13n\_1516\\13n\_1517\\13n\_1519\\13n\_1520\\13n\_1528\\13n\_1533\\13n\_1535\\13n\_1591\\13n\_1592\\13n\_1594\\13n\_1595\\13n\_1643\\13n\_1654\\13n\_1660\\13n\_1679\\13n\_1684\\13n\_1703\\13n\_1730\\13n\_1732\\13n\_1771\\13n\_1784\\13n\_1791\\13n\_1851\\13n\_1858\\13n\_1892\\13n\_1906\\13n\_1907\\13n\_1927\\13n\_1931\\13n\_1933\\13n\_1950\\13n\_1998\\13n\_2042\\13n\_2061\\13n\_2065\\13n\_2089\\13n\_2104\\13n\_2110\\13n\_2114\\13n\_2121\\13n\_2123\\13n\_2125\\13n\_2131\\13n\_2143\\13n\_2171\\13n\_2190\\13n\_2192\\13n\_2195\\13n\_2231\\13n\_2248\\13n\_2253\\13n\_2280\\13n\_2287\\13n\_2297\\13n\_2314\\13n\_2335\\13n\_2345\\13n\_2346\\13n\_2362\\13n\_2368\\13n\_2372\\13n\_2421\\13n\_2426\\13n\_2437\\13n\_2442\\13n\_2487\\13n\_2530\\13n\_2536\\13n\_2561\\13n\_2572\\13n\_2586\\13n\_2603\\13n\_2608\\13n\_2657\\13n\_2664\\13n\_2672\\13n\_2677\\13n\_2689\\13n\_2711\\13n\_2719\\13n\_2722\\13n\_2751\\13n\_2775\\13n\_2805\\13n\_2816\\13n\_2881\\13n\_2883\\13n\_2898\\13n\_2899\\13n\_2900\\13n\_2901\\13n\_2907\\13n\_2920\\13n\_2929\\13n\_2935\\13n\_2939\\13n\_2959\\13n\_2972\\13n\_2981\\13n\_2995\\13n\_3000\\13n\_3003\\13n\_3010\\13n\_3015\\13n\_3034\\13n\_3038\\13n\_3046\\13n\_3048\\13n\_3051\\13n\_3060\\13n\_3063\\13n\_3068\\13n\_3083\\13n\_3111\\13n\_3113\\13n\_3137\\13n\_3145\\13n\_3156\\13n\_3188\\13n\_3199\\13n\_3200\\13n\_3210\\13n\_3212\\13n\_3216\\13n\_3239\\13n\_3248\\13n\_3254\\13n\_3258\\13n\_3274\\13n\_3276\\13n\_3293\\13n\_3294\\13n\_3333\\13n\_3337\\13n\_3354\\13n\_3368\\13n\_3419\\13n\_3426\\13n\_3427\\13n\_3432\\13n\_3435\\13n\_3447\\13n\_3461\\13n\_3470\\13n\_3476\\13n\_3484\\13n\_3486\\13n\_3491\\13n\_3493\\13n\_3497\\13n\_3499\\13n\_3502\\13n\_3508\\13n\_3515\\13n\_3518\\13n\_3520\\13n\_3532\\13n\_3566\\13n\_3579\\13n\_3583\\13n\_3589\\13n\_3593\\13n\_3595\\13n\_3608\\13n\_3633\\13n\_3643\\13n\_3648\\13n\_3653\\13n\_3657\\13n\_3673\\13n\_3676\\13n\_3697\\13n\_3698\\13n\_3709\\13n\_3712\\13n\_3714\\13n\_3721\\13n\_3724\\13n\_3729\\13n\_3736\\13n\_3743\\13n\_3749\\13n\_3774\\13n\_3785\\13n\_3799\\13n\_3801\\13n\_3806\\13n\_3807\\13n\_3810\\13n\_3813\\13n\_3822\\13n\_3824\\13n\_3836\\13n\_3837\\13n\_3852\\13n\_3854\\13n\_3879\\13n\_3882\\13n\_3907\\13n\_3913\\13n\_3920\\13n\_3928\\13n\_3942\\13n\_3990\\13n\_4008\\13n\_4015\\13n\_4029\\13n\_4043\\13n\_4050\\13n\_4054\\13n\_4057\\13n\_4073\\13n\_4085\\13n\_4094\\13n\_4095\\13n\_4129\\13n\_4133\\13n\_4164\\13n\_4165\\13n\_4179\\13n\_4181\\13n\_4183\\13n\_4186\\13n\_4187\\13n\_4189\\13n\_4200\\13n\_4207\\13n\_4208\\13n\_4218\\13n\_4221\\13n\_4223\\13n\_4243\\13n\_4256\\13n\_4261\\13n\_4263\\13n\_4269\\13n\_4278\\13n\_4289\\13n\_4314\\13n\_4341\\13n\_4342\\13n\_4343\\13n\_4353\\13n\_4359\\13n\_4360\\13n\_4361\\13n\_4362\\13n\_4363\\13n\_4365\\13n\_4366\\13n\_4373\\13n\_4377\\13n\_4402\\13n\_4404\\13n\_4408\\13n\_4418\\13n\_4447\\13n\_4455\\13n\_4463\\13n\_4469\\13n\_4501\\13n\_4520\\13n\_4537\\13n\_4546\\13n\_4550\\13n\_4556\\13n\_4561\\13n\_4565\\13n\_4567\\13n\_4568\\13n\_4569\\13n\_4576\\13n\_4588\\13n\_4589\\13n\_4615\\13n\_4633\\13n\_4637\\13n\_4641\\13n\_4644\\13n\_4649\\13n\_4650\\13n\_4656\\13n\_4660\\13n\_4673\\13n\_4676\\13n\_4679\\13n\_4687\\13n\_4693\\13n\_4700\\13n\_4704\\13n\_4717\\13n\_4729\\13n\_4730\\13n\_4731\\13n\_4739\\13n\_4743\\13n\_4747\\13n\_4751\\13n\_4752\\13n\_4761\\13n\_4767\\13n\_4774\\13n\_4782\\13n\_4791\\13n\_4793\\13n\_4795\\13n\_4804\\13n\_4813\\13n\_4820\\13n\_4822\\13n\_4837\\13n\_4846\\13n\_4847\\13n\_4853\\13n\_4865\\13n\_4867\\13n\_4868\\13n\_4879\\13n\_4884\\13n\_4887\\13n\_4890\\13n\_4891\\13n\_4899\\13n\_4900\\13n\_4906\\13n\_4911\\13n\_4914\\13n\_4928\\13n\_4931\\13n\_4933\\13n\_4936\\13n\_4941\\13n\_4943\\13n\_4949\\13n\_4951\\13n\_4960\\13n\_4966\\13n\_4967\\13n\_4972\\13n\_4975\\13n\_4978\\13n\_4981\\13n\_4982\\13n\_4985\\13n\_4990\\13n\_4991\\13n\_4993\\13n\_4994\\13n\_4997\\13n\_5003\\13n\_5013\\13n\_5014\\13n\_5021\\13n\_5022\\13n\_5030\\13n\_5044\\13n\_5046\\13n\_5054\\13n\_5070\\13n\_5075\\13n\_5084\\13n\_5100\\13n\_5102}
\newcommand{\selftesttable}{\begin{tabular}{lr}
\toprule
corruption & rejected\\
\midrule
other smoothing of the crossing & 20 of 20\\
next crossing of the diagram & 13 of 20\\
wrong target ($12a_{1}$) & 20 of 20\\
bound lowered by one & 20 of 20\\
\bottomrule
\end{tabular}
}


\newtheorem{theorem}{Theorem}
\newtheorem{lemma}[theorem]{Lemma}
\newtheorem{corollary}[theorem]{Corollary}
\theoremstyle{definition}
\newtheorem{definition}[theorem]{Definition}
\theoremstyle{remark}
\newtheorem{remark}[theorem]{Remark}

\title[Pinch moves and the nonorientable four-genus]{Pinch moves and the nonorientable four-genus of knots with at most thirteen crossings}
\author{Deep Bhattacharjee}
\email{itsdeep@live.com}
\date{October 20, 2026}

\begin{document}
\begin{abstract}
The smooth nonorientable four-genus $\gamma_4(K)$ of a knot is the least first Betti number
of a smooth nonorientable surface in $B^4$ bounded by $K$. KnotInfo leaves it undetermined
for \ntotal{} prime knots with eleven to thirteen crossings. We determine it for
\nclosed{} of them and lower the known upper bound for \nimproved{} more. Each result comes
from one pinch move, the nonorientable band move at a crossing of a diagram, that turns the
knot into a slice knot or into a knot whose nonorientable four-genus is known. Every
certificate is a diagram, a crossing and a target knot. A program independent of the search
checks each one by Reidemeister moves or by an isomorphism of triangulated knot complements.
\end{abstract}
\maketitle

\section{Introduction}

Every knot $K\subset S^3=\partial B^4$ bounds a smooth surface in $B^4$; if $K$ is not
slice, it bounds nonorientable ones as well. The \emph{nonorientable four-genus} $\gamma_4(K)$ is the least
first Betti number of a smooth compact nonorientable surface $F\subset B^4$ with
$\partial F=K$; for slice knots one sets $\gamma_4(K)=1$. It was studied by Murakami and
Yasuhara \cite{MY} and by Gilmer and Livingston \cite{GiL}. Its values on small knots were computed by Jabuka and
Kelly for eight and nine crossings \cite{JK}, by Ghanbarian for ten \cite{Gh} and by Fairchild
for nonalternating knots with eleven \cite{Fa}. Lee and Sabloff \cite{LS} extended the
computation to fourteen crossings with an algorithmic search for M\"obius bands in the
style of Dunfield and Gong \cite{DG}, together with the Arf invariant obstruction of
Yasuhara \cite{Ya} and a linking form obstruction of Murakami and Yasuhara \cite{MY}. The KnotInfo database \cite{KnotInfo}
incorporates these data. In its version 2026.10.5, \ntotal{} knots with eleven to thirteen
crossings have an interval $[a,b]$ with $a<b$ as the entry for $\gamma_4$.

\begin{theorem}\label{thm:main}
For each of the \nclosed{} knots listed in the data file \texttt{certificates.json}
with the flag \texttt{closed}, $\gamma_4(K)$ equals the lower end of its KnotInfo interval:
\nclosedone{} knots have $\gamma_4(K)=1$ and \nclosedtwo{} have $\gamma_4(K)=2$. For the other
\nimproved{} knots of the file the upper bound is lowered by at least one.
\end{theorem}

Table~\ref{tab:counts} gives the counts by crossing number and by KnotInfo interval.
The proof uses one elementary move.

\begin{definition}
Let $D$ be a diagram of a knot $K$ and $c$ a crossing of $D$. Of the two smoothings of $c$,
one respects the orientation of $K$ and gives a two-component link; the other gives a knot
$K'$. We say $K'$ is obtained from $K$ by a \emph{pinch move} at $c$.
\end{definition}

\begin{lemma}\label{lem:pinch}
If $K'$ is obtained from $K$ by a pinch move, then $\gamma_4(K)\le 1$ when $K'$ is slice,
and $\gamma_4(K)\le\gamma_4(K')+1$ in general.
\end{lemma}

\begin{proof}
The pinch move is a band move, so there is a cobordism $\Sigma\subset S^3\times[0,1]$ from
$K$ to $K'$ formed by $K\times[0,1]$ and one band. Since $K'$ is connected, the band does not
respect the orientation of $K$, and $\Sigma$ is a M\"obius band with an open disk removed. If $K'$ bounds a disk $\Delta$ in $B^4$, then
$\Sigma\cup\Delta$ is a M\"obius band bounded by $K$, so $\gamma_4(K)\le1$. If $K'$ bounds
a surface $F$ with $b_1(F)=\gamma_4(K')$, then $\Sigma\cup F$ is a nonorientable surface
bounded by $K$ with $b_1=b_1(F)+1$. This is the bound used by Jabuka and Kelly \cite{JK}.
\end{proof}

The lower bounds are not ours. A knot that is not slice has $\gamma_4\ge1$. For the knots
whose interval starts at $2$, the lower bound is the obstruction to bounding a M\"obius band
recorded by KnotInfo from Lee and Sabloff \cite{LS}: either Yasuhara's condition
$4\operatorname{Arf}(K)-\sigma(K)\in\{0,\pm2\}\pmod 8$ fails \cite{Ya}, or the linking form
obstruction of Murakami and Yasuhara \cite{MY} applies. Theorem~\ref{thm:main} therefore rests on these
published obstructions, on the slice status and four-genus values recorded in KnotInfo for
the targets, and on the certificates below.

\section{Certificates and their verification}

A certificate for a knot $K$ consists of a planar diagram code $D$, a crossing $c$ of $D$, a
choice of smoothing, and a target: either the unknot or a knot $T$ of KnotInfo. It asserts
that
\begin{enumerate}
\item $D$ is a diagram of $K$;
\item the chosen smoothing of $c$ gives a one-component diagram $D'$;
\item $D'$ represents $T$ or its mirror image;
\item the bound $\beta=1$ if $T$ is slice or the unknot, and $\beta=\gamma_4(T)+1$ otherwise,
is smaller than the upper end of the KnotInfo interval of $K$.
\end{enumerate}
Since $\gamma_4$ and sliceness are invariant under mirror image, (1) to (4) and
Lemma~\ref{lem:pinch} give $\gamma_4(K)\le\beta$.

The certificates were found with SnapPy \cite{SnapPy}: for every open knot we applied
every pinch move to the KnotInfo diagram and to up to twenty-five diagrams obtained from it by
random Reidemeister moves, and identified the resulting knots by hyperbolic structure. A
second pass used the knots settled in the first pass as targets; the number of certificates
with such a target is \nchained{}, and the verifier checks them after the certificates of
their targets.

The verifier \texttt{verify.py} does not use SnapPy. It reads the KnotInfo database and
uses Regina \cite{Regina} for the topology. For (1) and (3) it proves that two diagrams
represent the same knot up to mirror image in one of two ways.
\begin{itemize}
\item \emph{Isotopy.} Both diagrams are simplified by Reidemeister moves to the same number
of crossings, and the closures of the two minimal diagrams under Reidemeister moves that add
at most one crossing share a diagram. Diagrams are compared by Regina's canonical
signatures, taken up to reflection, rotation and reversal.
\item \emph{Complements.} The triangulated complements of the two diagrams have
combinatorially isomorphic triangulations, possibly after Pachner moves that add at most one
tetrahedron. Then the complements are homeomorphic, and by the theorem of Gordon and Luecke
\cite{GL} the knots agree up to mirror image.
\end{itemize}
Both tests are exact: a positive answer is a proof. Regina's simplification is
randomised, so the verifier makes up to four attempts, each from a fixed seed, and gives each
certificate two minutes. For the unknot the verifier uses
Regina's recognition of the solid torus. Item (2) is checked by following the strands of
$D'$, and (4) by reading KnotInfo.

The verifier accepts all \ncerts{} certificates. It rejects certificates
corrupted in four ways, on a sample of twenty: the other smoothing of the crossing, the next
crossing, a wrong target, and a bound lowered by one (Table~\ref{tab:selftest}). The seven
modified certificates with the next crossing that it accepts are valid certificates: that
crossing gives the same target.

\section{Results}

\begin{table}[htbp]
\centering
\countstable
\caption{Open entries of $\gamma_4$ in KnotInfo 2026.10.5 and what the certificates give.
``Settled'': the value is now known; ``lowered'': the upper bound is lower.}
\label{tab:counts}
\end{table}

\begin{table}[htbp]
\centering
\selftesttable
\caption{Corrupted certificates rejected by the verifier.}
\label{tab:selftest}
\end{table}

The remaining \nopen{} knots, listed in \texttt{still\_open.json}, are not settled by one
pinch move found in this way. For a knot with interval $[1,2]$ that is not settled, no
pinch move we tried gives a slice knot; this says nothing about whether $\gamma_4=1$.

\begin{remark}
The \nclosedone{} M\"obius bands found here are not in the data of Lee and Sabloff
\cite{LS}, whose search follows Dunfield and Gong \cite{DG} and runs over a finite family of
bands chosen in advance. We have not determined which restriction of that family excludes
them. The knots settled at $\gamma_4=2$ were outside their aim, which was M\"obius bands;
for them the target of the pinch move is a knot known to bound a M\"obius band but not a disk.
\end{remark}

\appendix
\section{The settled knots}\label{app:list}

{\footnotesize
\noindent\textbf{$\gamma_4=1$ (\nclosedone{} knots).}
\begin{multicols}{6}\raggedright
\listone
\end{multicols}

\noindent\textbf{$\gamma_4=2$ (\nclosedtwo{} knots).}
\begin{multicols}{6}\raggedright
\listtwo
\end{multicols}
}

\begin{thebibliography}{99}
\bibitem{DG} N.~M.~Dunfield and S.~Gong, \emph{Ribbon concordances and slice obstructions:
experiments and examples}, arXiv:2512.21825.
\bibitem{Fa} M.~Fairchild, \emph{The non-orientable 4-genus of 11 crossing non-alternating
knots}, J. Knot Theory Ramifications 34 (2025), no.~1.
\bibitem{Gh} N.~Ghanbarian, \emph{The non-orientable 4-genus for knots with 10 crossings},
J. Knot Theory Ramifications 31 (2022), no.~5, 2250034.
\bibitem{GiL} P.~M.~Gilmer and C.~Livingston, \emph{The nonorientable 4-genus of knots},
J. Lond. Math. Soc. (2) 84 (2011), 559--577.
\bibitem{GL} C.~McA.~Gordon and J.~Luecke, \emph{Knots are determined by their complements},
J. Amer. Math. Soc. 2 (1989), 371--415.
\bibitem{JK} S.~Jabuka and T.~Kelly, \emph{The nonorientable 4-genus for knots with 8 or 9
crossings}, Algebr. Geom. Topol. 18 (2018), 1823--1856.
\bibitem{KnotInfo} C.~Livingston and A.~H.~Moore, \emph{KnotInfo: Table of Knot Invariants},
knotinfo.org, version 2026.10.5.
\bibitem{LS} D.~Lee and J.~M.~Sabloff, \emph{An algorithmic search for knots bounding
M\"obius bands in $B^4$}, arXiv:2607.23582.
\bibitem{MY} H.~Murakami and A.~Yasuhara, \emph{Four-genus and four-dimensional clasp number
of a knot}, Proc. Amer. Math. Soc. 128 (2000), 3693--3699.
\bibitem{Regina} B.~A.~Burton, R.~Budney, W.~Pettersson et al., \emph{Regina: software for
low-dimensional topology}, regina-normal.github.io, version 7.4.
\bibitem{SnapPy} M.~Culler, N.~M.~Dunfield, M.~Goerner and J.~R.~Weeks, \emph{SnapPy, a
computer program for studying the geometry and topology of 3-manifolds}, snappy.computop.org.
\bibitem{Ya} A.~Yasuhara, \emph{Connecting lemmas and representing homology classes of simply
connected 4-manifolds}, Tokyo J. Math. 19 (1996), 245--261.
\end{thebibliography}

\end{document}
